\textbf{Pour cette première partie, aucune justification n'est demandée et une seule réponse est possible par question.} Les justifications ci-dessous sont fournies à titre pédagogique.

Une réponse fausse ou l'absence de réponse n'enlève aucun point : il ne faut donc jamais laisser une question sans réponse.

\subsection*{Question 1 \quad Réponse A}

On développe l'opposé $-(a-b) = -a + b$, puis on réduit :
\[E = -(a-b) + a - b = -a + b + a - b = -a+a + b-b = 0.\]

\subsection*{Question 2 \quad Réponse D}

$x^2 = 6$ admet deux solutions opposées :
\[x^2 = 6 \iff x = \sqrt{6} \ \text{ ou } \ x = -\sqrt{6}.\]
D'où $S = \left\{-\sqrt{6}~;~\sqrt{6}\right\}$.

\subsection*{Question 3 \quad Réponse D}

On applique l'identité remarquable $(a-b)^2 = a^2 - 2ab + b^2$ avec $a = 3x$ et $b = 5$ :
\[(3x-5)^2 = (3x)^2 - 2\times 3x \times 5 + 5^2 = 9x^2 - 30x + 25.\]

\subsection*{Question 4 \quad Réponse A}

La droite coupe l'axe des ordonnées en 3 : l'ordonnée à l'origine est $b = 3$. Elle passe
aussi par les points de coordonnées $(0~;~3)$ et $(6~;~0)$ donc son coefficient directeur vaut :
\[a = \dfrac{0 - 3}{6 - 0} = -\dfrac{1}{2}.\]
Une équation de la droite est donc $y = -\dfrac{1}{2}x + 3$.

\subsection*{Question 5 \quad Réponse C}

On applique la règle $\frac{a^n}{a^m}=a^{n-m}$ :
\[\dfrac{2^{10}}{2} = \dfrac{2^{10}}{2^1} = 2^{10-1} = 2^9.\]

\subsection*{Question 6 \quad Réponse B}

On isole $x^2$ puis on prend la racine carrée (les trois réels sont strictement positifs,
donc $x > 0$) :
\[a = \dfrac{1}{2}bx^2 \iff 2a = bx^2 \iff x^2 = \dfrac{2a}{b} \iff x = \sqrt{\dfrac{2a}{b}}.\]

\subsection*{Question 7 \quad Réponse C}

$f(x) = (2x-6)(-x+1)$ est un produit de deux facteurs affines, qui s'annulent respectivement
en 3 et en 1. Le coefficient de $x^2$ dans le développement de $f(x)$ vaut
$2 \times (-1) = -2 < 0$ : la parabole associée est ouverte vers le bas, donc $f$ est
négative sur $]-\infty ; 1[$, positive sur $]1 ; 3[$ et négative sur $]3 ; +\infty[$.

\subsection*{Question 8 \quad Réponse A}

Une baisse de 20\,\% correspond au coefficient multiplicateur $0,8$ ; une baisse de
30\,\%, au coefficient multiplicateur $0,7$. Le coefficient global est le produit
des deux :
\[0,8 \times 0,7 = 0,56 = 1 - 0,44.\]
Les deux baisses successives sont donc équivalentes à une baisse globale de 44\,\%.
